\subsection{Pérdida de Wasserstein}

Hu et al.~\cite{hu2019} propusieron la primera pérdida basada en
persistencia. Se calcula $\operatorname{Dgm}(L)$ con la filtración de
supernivel, se empareja con $\operatorname{Dgm}(G)$ mediante el
emparejamiento óptimo $\gamma^*$ de Wasserstein y se minimiza
$\sum_{q}\|q-\gamma^*(q)\|_2^2$. Como $G$ es binaria, todos los puntos de
$\operatorname{Dgm}(G)$ coinciden en el punto $(0,1)$. La
pérdida empuja hacia $(0,1)$ las características emparejadas y hacia la
diagonal las demás. El problema es el señalado en la
Sección~\ref{sec:distancias}: $\gamma^*$ puede emparejar un agujero de la
predicción con otro de la referencia situado en una zona distinta de la
imagen~\cite{stucki2023}.

Por ejemplo, si la predicción tiene un agujero cerca de la esquina superior
izquierda y la referencia uno cerca de la inferior derecha, ambos con la misma
persistencia, el emparejamiento de Wasserstein puede unirlos aunque no
correspondan al mismo objeto: el diagrama no contiene ninguna información
sobre la posición.

\subsection{Betti matching}
\label{sec:betti-matching}

Stucki et al.~\cite{stucki2023} resuelven este problema construyendo un
emparejamiento que respeta la posición espacial. La herramienta son los
\emph{emparejamientos inducidos}.

\begin{definicion}[{\cite{bauer2015}}]
Un \emph{morfismo} de módulos de persistencia $\Phi:M\to N$ es una familia de
aplicaciones lineales $\Phi_t:M_t\to N_t$, una por cada valor $t$ del
parámetro, que conmuta con las aplicaciones de transición:
$N_{s\le t}\circ\Phi_s=\Phi_t\circ M_{s\le t}$ para todo $s\le t$. Su
\emph{módulo imagen} es $\operatorname{im}\Phi$, con
$(\operatorname{im}\Phi)_t=\operatorname{im}\Phi_t\subseteq N_t$ y las
transiciones restringidas de $N$.
\end{definicion}

El módulo imagen es un submódulo de $N$: contiene las clases de $M$ que
sobreviven a través de $\Phi$. Es precisamente lo que relaciona el teorema
siguiente con los códigos de barras de $M$ y $N$. Sus intervalos aparecen como
intervalos de $M$ y de $N$ a la vez, en el sentido del teorema siguiente.

\paragraph{Emparejamientos inducidos.} Un morfismo $\Phi:M\to N$ entre
módulos de persistencia determina un emparejamiento entre sus códigos de
barras. Se trata de un caso particular del teorema 4.2 de~\cite{bauer2015}:

\begin{teorema}[{\cite[Teor.~4.2]{bauer2015}, \cite[Teor.~1]{stucki2023}}]
Sea $\Phi:M\to N$ un morfismo de módulos de persistencia (de dimensión
finita puntual, escalonados y continuos por arriba). Existen inyecciones
únicas $\mathcal{B}(\operatorname{im}\Phi)\hookrightarrow\mathcal{B}(M)$ y
$\mathcal{B}(\operatorname{im}\Phi)\hookrightarrow\mathcal{B}(N)$ que envían
cada intervalo $[b,c)$ a un intervalo $[b,d)\in\mathcal{B}(M)$ con $c\le d$
y a un intervalo $[a,c)\in\mathcal{B}(N)$ con $a\le b$.
\end{teorema}

Componiendo ambas inyecciones a través de la imagen se obtiene el
\emph{emparejamiento inducido} $\sigma(\Phi):\mathcal{B}(M)\to\mathcal{B}(N)$.
Para imágenes, si $I\ge J$ píxel a píxel, cada subnivel de $I$ está
contenido en el de $J$. Las inclusiones inducen un morfismo
$\Phi(I,J):H_*(I)\to H_*(J)$ y, por tanto, un emparejamiento
$\sigma(I,J)$. Este emparejamiento identifica características
\emph{a través de la inclusión}, es decir, en la misma posición.

\paragraph{Imagen de comparación~\cite[§3.1]{stucki2023}.} Para comparar $L$ y $G$, que no son
comparables entre sí, se introduce un espacio ambiente común. Con
supernivel se toma $C=\max(L,G)$, cuyos supernivel son la unión de los de
$L$ y $G$. Como $L,G\le C$, hay emparejamientos inducidos $\sigma(L,C)$ y
$\sigma(G,C)$, y el \emph{Betti matching} es la composición
\[
  \tau(L,G)=\sigma(G,C)^{-1}\circ\sigma(L,C).
\]
Dos características quedan emparejadas si ambas sobreviven en la unión como
la misma característica.

\paragraph{Pérdida~\cite[§3.2]{stucki2023}.} Viendo los códigos de barras como diagramas, y
emparejando con la diagonal los puntos no emparejados, la pérdida de Betti
matching es
\[
  \ell_{BM}(L,G)=\sum_{q\in\operatorname{Dgm}(L)}2\,\|q-\tau(L,G)(q)\|_2^2,
\]
y en la práctica se combina con una pérdida de solapamiento,
$\ell=\alpha\,\ell_{BM}+\ell_{\mathrm{Dice}}$. Cada punto $q$ depende
solo de dos píxeles, el de nacimiento $b(q)$ y el de muerte $d(q)$. Por
eso el gradiente es disperso: refuerza el contraste entre esos dos píxeles
para las características emparejadas y lo reduce para las no emparejadas.

\paragraph{Métrica~\cite[§3.3]{stucki2023}.} Para imágenes binarias todos los puntos están en
$(0,1)$. Un punto no emparejado contribuye
$2(\tfrac12)^2+2(\tfrac12)^2=1$ y uno emparejado contribuye $0$, de modo que
el \emph{error de Betti matching} es un entero:
\[
  \tau^{\mathrm{err}}(P,G)=\#\ker\tau(P,G)+\#\operatorname{coker}\tau(P,G),
\]
el número de características de $P$ o de $G$ sin pareja espacialmente
correcta. Refina el error de Betti, que se obtiene con el mismo recuento
pero usando un emparejamiento que ignora la posición.